\begin{frame}[t, fragile]{Type classes in Haskell.}

We could then create an instance of our new type as follows;

\begin{minted}[linenos=true]{haskell}
data Person = Person String Int
    deriving Show

main :: IO ()
main = do

    let alice = Person "Alice" 30

    print alice
\end{minted}

Compiling and running this program will produce the following output;

\begin{minted}[linenos=true]{console}
Person "Alice" 30
\end{minted}

If we don't like the look of this output, then we can rewrite our \mintinline[fontsize=\normalsize]{haskell}{Person}
type such that it defines its own \mintinline[fontsize=\normalsize]{haskell}{show} function. This will be shown in the
next section.

\end{frame}


% New slide.

\begin{frame}[t, fragile]{Implementing a custom version of the \mintinline[fontsize=\normalsize]{haskell}{show} function.}

\begin{minted}[linenos=true]{haskell}
instance Show Person where
    show (Person name age) =
        "Person: " ++ name ++ ", age " ++ show age
\end{minted}

Note in this code snippet, the presence of \mintinline[fontsize=\normalsize]{haskell}{instance Show} rather than
\mintinline[fontsize=\normalsize]{haskell}{deriving Show} as was used in an earlier code snippet.

\end{frame}


\begin{frame}[t, fragile]{Instantiating types in Haskell.}

You can't create variables in Haskell.

However you can create what Haskell calls ``bindings''.

Consider the following Haskell code which simply creates a binding named \mintinline[fontsize=\normalsize]{haskell}{my_var_1} and then prints its value to the
screen.

\begin{minted}[linenos=true]{haskell}
import Text.Printf (printf)

main :: IO ()
main = do

    let

        my_var_1 :: Int
        my_var_1  = 1

    printf "my_var_1 = %d\n" my_var_1
\end{minted}

When it comes to creating the binding \mintinline[fontsize=\normalsize]{haskell}{my_var_1}, the lines of code which are relevant, are the two lines of code
which follow the let statement. The first line of code declares a binding called \mintinline[fontsize=\normalsize]{haskell}{my_var_1} and declares it to be of
the standard Haskell data type Int. The next line of code then initialises \mintinline[fontsize=\normalsize]{haskell}{my_var_1} to a value of 1.

\end{frame}